\begin{proof}
\begin{enumerate}[label=\arabic*.]
\item We first identify the actual moving pencil on each known member.
By \cref{night:P1A.1}, $K_M^{-4}=h^*L^{24}$ and the fibres of $h$
are connected.  For the coarse map
$\rho:\widehat B\to\mathbb P^1$, the root-divisor formula gives
\[
 \rho_*L^{24}
 =\mathcal O_{\mathbb P^1}
   \bigl(\lfloor96/5\rfloor\,\infty+
                \lfloor-72/4\rfloor\,0\bigr)
 =\mathcal O_{\mathbb P^1}(19\infty-18\,0).
\]
Its sections are spanned by $\lambda^{18},\lambda^{19}$.
On a root chart the invariant sections have a common local
monomial factor, and their ratio cancels that factor.  Thus
$h^0(M,K_M^{-4})=2$, and the moving part of this complete
system is exactly the marked base map, up to a projective
coordinate change.  In particular its base ideal is invertible
on every member of the prescribed family.

\item Suppose that a proper Hausdorff analytic extension over a full
parameter disc exists.  Take the reduced component containing
the given nonempty open family and resolve it.  Analytic
resolution applies on a relatively compact neighbourhood of
the compact central fibre and can be kept unchanged over the
given smooth locus
\cite[Theorem 1.6 and its following analytic-neighbourhood
paragraph, preprint p.~5]{BM97}.
Write $p:Z\to\Delta$ for the resulting proper map from a smooth
irreducible space.  The proper images of the exceptional and
critical loci are analytic
\cite[\S7, Satz 1, p.~60]{Grauert60}; they omit the given open
parameter set.  After shrinking the disc they are supported
at its centre.  Hence $p$ is a proper holomorphic submersion
over $\Delta^*$, with compact fourfold fibres.

Put $\mathcal J=\omega_{Z/\Delta}^{-4}$.
The sheaves $R^jp_*\mathcal J$ are coherent
\cite[\S6, Hauptsatz I, p.~59]{Grauert60}.
Multiplication by a nonzero parameter germ is injective on
$\mathcal J$ and on $p_*\mathcal J$, so $p_*\mathcal J$ is locally free
over the disc and free after shrinking.
The actual exact sequence for a fibre over $t$ gives
\[
 0\longrightarrow(p_*\mathcal J)_t/\mathfrak m_t(p_*\mathcal J)_t
 \longrightarrow H^0(Z_t,\mathcal J_t)
 \longrightarrow
 \ker\bigl((t'-t): (R^1p_*\mathcal J)_t
                       \longrightarrow(R^1p_*\mathcal J)_t\bigr)
 \longrightarrow0.
\]
The last term is supported in the torsion support of the
coherent sheaf $R^1p_*\mathcal J$, a closed discrete analytic
set.  At a point of the given open family outside this support,
the displayed comparison shows that $p_*\mathcal J$ has rank
two.  Shrinking makes the support empty away from the centre.
Thus a frame of $p_*\mathcal J$ gives the complete two-dimensional
linear system on each punctured fibre.  This uses the displayed
coherent maps, without assuming closedness of a general
base-change failure set.

\item The two evaluation sections define a coherent base ideal and
a meromorphic map $Z\dashrightarrow\mathbb P^1$ over the full
disc.  The ideal is locally invertible on the given open family
by the first step.  Its non-invertibility locus is analytic;
its proper image in $\Delta$ is therefore discrete and can be
removed away from the centre.  Blowing up that ideal and
resolving the graph consequently changes only the central
fibre.  It gives a proper holomorphic map to
$\mathbb P^1\times\Delta$.  The construction uses the graph
of the coherent evaluation map, so its coefficients have
finite meromorphic order in the parameter.

Its critical image is proper analytic.  Every component
dominating the parameter disc meets the given open family,
where there are precisely three critical values.  Their
degrees over the parameter are one.  The multiplicities five
and four distinguish two of them, and the third is the cusp.
After shrinking, their normalizations give three distinct
labelled meromorphic sections, and all vertical exceptional
components lie over the centre.  The projective transformation
taking these sections to $\infty,0,\lambda_*$ has meromorphic
coefficients and is invertible on $\Delta^*$.  Resolving its
graph gives the actual marked map $H$ over the disc, still
unchanged on punctured fibres.  Its regular restriction is a
proper holomorphic submersion
\[
 q:W\longrightarrow
 B=\bigl(\mathbb P^1\setminus\{0,\infty,\lambda_*\}\bigr)
          \times\Delta^*.
\]
Ehresmann transport makes the number of components of a fibre
locally constant.  The base is connected and the known fibres
are connected, so every fibre of $q$ is connected.

\item We need the torus structure on an analytically open dense
part of this regular family.  The coherent sheaf
$\mathcal E=q_*\Omega^1_{W/B}$ has generic rank three.
Generic cohomology and base change for this flat bundle gives
a proper analytic exceptional set outside which $\mathcal E$
is locally free of rank three and its fibre is the space of
holomorphic one-forms.  The constant-dimension comparison is
\cite[\S7, Satz 5, pp.~63--64]{Grauert60}.
Use the global coherent evaluation map
\[
 q^*\mathcal E\longrightarrow\Omega^1_{W/B}.
\]
Its cokernel has analytic support.  The proper image of that
support is analytic in $B$ and omits the nonempty part of the
known torus open set outside the generic exceptional set.
Enlarge the exceptional set by this image.  On its complement
evaluation is an isomorphism of rank-three bundles.

Relative exterior differentiation defines an
$\mathcal O_B$-linear map
$\mathcal E\to q_*\Omega^2_{W/B}$.  It vanishes on the known
torus open set.  The target is torsion-free, since it is the
proper direct image of a vector bundle under a submersion.
The identity theorem therefore makes this map zero.  Each
remaining compact connected fibre has three everywhere
independent closed holomorphic one-forms.  Their dual vector
fields commute and are complete by compactness.  Their flows
give a transitive holomorphic $\mathbb C^3$-action: every orbit
is open and the fibre is connected.  Its stabilizer $\Gamma$
is discrete, and compactness of $\mathbb C^3/\Gamma$ makes
$\Gamma$ a full real lattice.  This is the actual complex-torus
structure on each such fibre.

\item Let $E\subset B$ be the proper analytic exceptional set just
removed.  On the universal cover of $B\setminus E$, use
Gauss--Manin transport from the given open family to mark
integral $H^1$.  Integration of the closed relative forms
gives a holomorphic map to $\operatorname{Gr}(3,6)$.  On the
known open set it is the explicit graph
$F_\xi=(\sigma+\xi\alpha e)(P)$.  The identity theorem on the
connected cover identifies these two maps everywhere.
The explicit graph is a rank-three holomorphic plane for
every finite $\xi$, including parameters for which it is not
real-transverse.  At infinity the same assertion is used with
$\xi=1/u$ on the punctured parameter disc.  All changes of
continuation marking are real integral and preserve the
condition $F\cap\bar F=0$.

\item Let $\xi_0$ be a finite boundary point, so
$\operatorname{Im}\xi_0=c_*=-\min h$.
The minimum may occur at a finite root point.  Continuity
there nevertheless gives regular $\lambda$ with $h(\lambda)$
arbitrarily close to $\min h$.  Every sufficiently small full
parameter disc about $\xi_0$ therefore contains noncentral
points of
\[
 \operatorname{Im}\xi+h(\lambda)=0.
\]
This equation defines nonempty smooth real hypersurface
pieces in the complex surface $B$, since its derivative in
$\operatorname{Im}\xi$ is one.  A proper complex analytic
subset of $B$ has real dimension at most two; it cannot
contain those three-dimensional pieces.  Choose such a
point outside $E$.  Its fibre is a complex torus by the fourth
step, so its period plane is real-transverse.  The fifth step
identifies that plane with $F_\xi$, whereas the exact criterion
of \cref{night:P1.8} gives
$F_\xi\cap\overline{F_\xi}\ne0$ at the chosen point.
This is a contradiction.

At infinity fix a regular $\lambda$ and set
\[
 \xi=R-i h(\lambda),\qquad
 u=\frac1{R-i h(\lambda)},\qquad R\in\mathbb R.
\]
These are noncentral points of the same failure locus and
tend to $u=0$ as $|R|\to\infty$.  Varying $\lambda$ and $R$
again gives real hypersurface pieces in every punctured
$u$-disc, so one can choose a point outside $E$ and obtain the
same contradiction.  Under a finite ramified change of the
boundary coordinate, the parameter map has nonzero derivative
off the centre after shrinking.  The inverse images of these
hypersurface pieces remain nonempty real hypersurfaces there,
and the argument is unchanged.

All contradictions occur on a noncentral compact smooth
torus fibre of the recovered marked map.  The reductions and
modifications are proper and preserve the prescribed open
family; they impose no smoothness, multiplicity or component
condition on the original central fibre.  The permitted
central spaces therefore do not avoid the contradiction.
Both nonextension assertions follow.
\end{enumerate}
\end{proof}
